\section*{4. Griffiths 定理}

\subsubsection*{定理 4.1.}
设 $X$ 在 $k$ 上射影、光滑且连通, 偶数维数为 $n=2m\geq4$. 设 $D$ 为由次数为 $d$ 的超曲面截面组成的 Lefschetz 线束. 假设:

\medskip
\noindent\textbf{(4.1.0)} $X$ 满足 (LV) (XVIII~5.2.2);

\noindent\textbf{(4.1.1)} 束 $D$ 满足假设 (A) (XVIII~5.3);

\noindent\textbf{(4.1.2)} 束 $D$ 满足另一种情形 (3.2.2), 即在命题~3.2 的设定下, 消灭上同调 $E^{n-1}(X_{\bar\eta},\mathbb Q_\ell)$ 的水平并非 $\leq1$;

\noindent\textbf{(4.1.3)} 包含映射 $m:X_{\bar\eta}\hookrightarrow X\times_k k(\bar\eta)$ 满足假设 (B) (3.3).\footnote{印本在这个几何包含映射前写作 $m^*$. 3.3 中的条件 (B) 是对 $m$ 施加的; $m^*$ 表示它所诱导的上同调映射.}

\medskip
则对于 $X$ 上余维为 $m$ 的代数圈 $z$, 若其在 $X_{\bar\eta}$ 上的限制作为代数圈满足
\begin{equation}\tag{4.1.4}
z_{\bar\eta}=z|_{X_{\bar\eta}}
\quad\text{代数等价于零},
\end{equation}
则有
\begin{equation}\tag{4.1.5}
\operatorname{cl}(z)=0
\quad\text{于 }H^{2m}(X,\mathbb Q_\ell(m));
\end{equation}
换言之, $z$ 在 $\mathbb Q_\ell$-同调中等价于零.

\subsubsection*{4.2. 关于 $X$ 为完全交的情形的一点说明}
如前所述 (参见~XVIII~5.2.5 和 (3.4)), 此时假设 (4.1.0) 和 (4.1.3) 自动成立. 在 XVIII~6.3.4 和 6.4 中已见, 次数为 $d>2n$ 的超曲面束满足假设 (4.1.1). 此外, 下一讲中将对次数为 $d>2n$ 的``一般''超曲面束证明条件 (4.1.2) (XXI~5.2). 因此, 次数为 $d>2n$ 的``一般''超曲面束满足定理~4.1 的全部假设.

\subsubsection*{4.3.}
回到任意 $X$ 的情形, 很有理由相信: 当 $d$ 充分大时, 至少对于充分一般的 $D$, 4.1 的全部条件均对 $D$ 成立. 对于 (4.1.1) 已经看到这一点 (XVIII~6.4), 而且已提到 (4.1.0) 和 (4.1.3) 总成立的猜想 (当 $\operatorname{char}(k)=0$ 时, (4.1.0) 在一切情形下均已知成立). 至于条件 (4.1.2), 当 $\operatorname{char}(k)=0$ 时, Hodge 理论证明其对充分大的 $d$ 成立,\footnote{印本在这个有关 Hodge 理论的分句中写作 $\operatorname{char}(k)\ne0$. 随后的分句处理正特征情形, 故此处应读作 $\operatorname{char}(k)=0$.} 而下一讲附录中所给的启发性论证表明, 当 $\operatorname{char}(k)>0$ 时, 同一结论也很可能成立.

更一般地, 这些论证---分别为 Hodge 理论和晶体上同调---证明了或分别使人相信: 对于维数为 $n$ 的光滑射影簇 $X$, 对每个次数 $d$ 充分大的超曲面的 Lefschetz 线束 $D$ (分别地, 对每个次数为 $d$ 且充分一般的超曲面束 $D$), 消灭上同调 $E^{n-1}(X_{\bar\eta},\mathbb Q_\ell)$ 的水平恰为 $n-1$; 即其水平并非 $\leq n-3$, 或等价地, 其余水平并非 $1$.
